% !TeX root = ../../Book3.tex
% 纲要标题：### C.3 余切复形、切复形与形变问题

\section{余切复形、切复形与形变问题}
\label{b3:sec:C03}

% 纲要细目（与计划纲要同步）：
% * 余切复形（cotangent complex）与切复形（tangent complex）及形变理论的基本问题；原有切复形主题保留，但与 $L_{B/A}$ 及其在适当假设下的导出对偶分别命名
%
% <!-- 短节：不设小节 -->


% 正文以具体模型说明专题；长证明给出概要并指明技术细节的出处。

Kähler 微分通过 \(\mathrm{d}f\) 记录方程的一阶变化，但它是一个商模，已把产生关系的映射压缩掉了。形变问题往往还需要知道关系怎样变化，以及关系之间是否有进一步关系。\chapref{b3:ch:23}的平方零提升计算已经出现了这种需求。

\ProseParagraphBreak
先固定可计算的范围：设 \(k\) 为特征零域，\(P=k[x_1,\ldots,x_n]\)，\(I=(f_1,\ldots,f_c)\) 由正则序列生成，\(B=P/I\)。完全交情形的\term{余切复形}[cotangent complex]有两项模型
\[
L_{B/k}=\bigl[I/I^2\xrightarrow{\mathrm{d}}\Omega_{P/k}\otimes_PB\bigr]
\simeq\bigl[B^c\xrightarrow{J}B^n\bigr],
\]
次数分别为 \(-1,0\)，其中 \(J(e_j)=\sum_i(\partial f_j/\partial x_i)\,\mathrm{d}x_i\)。\(H^0(L_{B/k})=\Omega_{B/k}\) 由余法正合列得到。一般非完全交不能直接把这两项当作全部余切复形；还须用更长的消解记录高阶关系。下面说明完全交为什么恰好得到这两项；一般环同态的无界构造仍须另作系统学习。

\begin{proofsketch}
使用自由微分分次代数计算余切复形：给 \(P\) 添加次数 \(-1\) 的反交换生成元 \(\xi_1,\ldots,\xi_c\)，令
\[
Q=P\otimes_k\bigwedge(\xi_1,\ldots,\xi_c),
\qquad \mathrm{d}\xi_j=f_j,\qquad \mathrm{d} x_i=0.
\]
正则序列的 Koszul 正合性，见\cref{b3:prop:ci-koszul-resolution}，说明 \(Q\to B\) 是一个自由微分分次代数消解。定义上，余切复形由
\(B\otimes_Q\Omega_{Q/k}\) 计算。微分模自由地由次数零的 \(\mathrm{d}x_i\) 和次数 \(-1\) 的 \(\mathrm{d}\xi_j\) 生成；复形微分把后者送到 \(\mathrm{d}f_j\)。张量到 \(B\) 后便恰为 \([B^c\xrightarrow{J}B^n]\)。正则序列还使 \(I/I^2\) 以各 \(f_j\) 的类为基，给出所写的 \(I/I^2\) 形式。

这一计算没有依赖展示的数学内容在于：不同自由微分分次代数消解之间，可以按自由生成元逐阶提升同态，并同样逐阶构造两种提升间的同伦；所得微分模在张量到 \(B\) 后因而给出同一个拟同构类。本概要省略这个一般消解比较引理的归纳，保留这里所需的有限模型。低阶余切理论、完全交与形变解释可对照\cite[Sections 2.1, 3.2 and 5.1]{LichtenbaumSchlessinger1967CotangentComplex}。正则序列假设在 \(Q\to B\) 的正合性中起作用；去掉它就不能使用这个两项计算。
\end{proofsketch}

\ProseParagraphBreak
上述模型的两项都是有限自由模，所以逐项取对偶已计算其导出对偶。所得
\[
T_{B/k}=\operatorname{Hom}_B(L_{B/k},B)
=\bigl[B^n\xrightarrow{J^{\mathsf T}}B^c\bigr]
\]
位于次数 \(0,1\)，称为\term{切复形}[tangent complex]的模型。零次核就是保持全部关系的导子 \(\operatorname{Der}_k(B,B)\)，一次余核则度量不能由坐标变化消掉的方程扰动。余切复形与切复形因此是对偶关系，不是同一记号的两个名称。一般非自由情形须先作消解，再用导出 Hom，不能随意逐项对偶。

\begin{definition}\label{b3:def:first-order-deformation}
设 \(k\) 为域，\(B\) 为交换 \(k\) 代数，\(D=k[\varepsilon]/(\varepsilon^2)\)。若平坦的交换 \(D\) 代数 \(B_\varepsilon\) 连同指定同构 \(B_\varepsilon/\varepsilon B_\varepsilon\cong B\) 一起给出，则称这组数据为 \(B\) 的\term{一阶形变}[first-order deformation]。形变之间的同构须与该指定同构相容。
\end{definition}

平坦性排除了任意添加幂零关系。例如 \(D[x]/(\varepsilon x)\) 在模 \(\varepsilon\) 后也是 \(k[x]\)，却把 \(x\) 在无穷小方向额外零化，因而不是所需的平坦变化。一般定义及其上同调解释见\cite[Section 4.3, Definition 4.3.2 and Theorem 4.3.3]{LichtenbaumSchlessinger1967CotangentComplex}。下例不依赖这个一般定理，直接计算方程扰动。

\begin{proposition}[超曲面的方程扰动]\label{b3:prop:hypersurface-first-order-perturbation}
设 \(k\) 为特征零域，\(P=k[x_1,\ldots,x_n]\)，\(0\ne f\in P\)，\(B=P/(f)\)，\(D=k[\varepsilon]/(\varepsilon^2)\)。对任意 \(g\in P\)，代数
\[
B_g=D[x_1,\ldots,x_n]/(f+\varepsilon g)
\]
为 \(B\) 的一阶形变。若把扰动按变换 \(x_i\mapsto x_i+\varepsilon a_i\) 及生成方程乘 \(1+\varepsilon b\) 视为等价，则等价类由
\[
P/(f,\partial f/\partial x_1,\ldots,\partial f/\partial x_n)
\]
参数化。令 \(T_{B/k}=[B^n\xrightarrow{(\partial f/\partial x_i)}B]\) 为次数 \(0,1\) 上的切复形模型，则这个参数模也同构于 \(H^1(T_{B/k})\)。
\end{proposition}
\begin{proof}
选取 \(k\) 向量空间补空间 \(V\)，使 \(P=V\oplus fP\)。将 \(h_0+\varepsilon h_1\) 中 \(h_0\) 写成 \(v_0+fq_0\)，再用 \(f=-\varepsilon g\) 代换，并将 \(h_1-gq_0\) 模 \(fP\) 化简，得到唯一形式 \(v_0+\varepsilon v_1\)，其中 \(v_i\in V\)。唯一性如下：若该式等于 \((f+\varepsilon g)(q_0+\varepsilon q_1)\)，比较常数项得 \(v_0=fq_0\)，所以 \(v_0=q_0=0\)；再比较 \(\varepsilon\) 项得 \(v_1=fq_1\)，故 \(v_1=0\)。这里用到 \(P\) 是整环且 \(f\ne0\)。因此 \(B_g\) 作为 \(D\) 模自由，同构于 \(D\otimes_kV\)，故平坦。

平方零 Taylor 公式给出
\[
f(x+\varepsilon a)+\varepsilon g(x+\varepsilon a)
=f+\varepsilon\Bigl(g+\sum_i a_i\frac{\partial f}{\partial x_i}\Bigr).
\]
再乘 \(1+\varepsilon b\) 会给扰动加上 \(bf\)。这些变换都有把 \(\varepsilon\) 项取相反数得到的逆，且其复合仍只对 \(g\) 添加上述理想中的元素。因此等价关系恰是模该理想同余。最后在 \(B\) 中取 Jacobian 行向量的余核，得到同一个商。
\end{proof}

\begin{example}\label{b3:exm:cusp-tangent-complex}
对尖点 \(B=k[x,y]/(y^2-x^3)\)，\(\operatorname{char}k=0\)，上式给出
\[
H^1(T_{B/k})=B/(x^2,y)\cong k[x]/(x^2).
\]
所以任意方程扰动在上述等价下可化为 \(y^2-x^3+\varepsilon(a+bx)\)，且两个参数不能被这种坐标变化消去。若只检查光滑点的切空间，便看不到这一整族方程的变化。
\end{example}

从 \(\varepsilon^2=0\) 延长到更高次幂零参数时，还须检查新的相容性；一般形变的障碍由更高的上同调群表达。这里对超曲面的一阶计算说明了这种理论的对象，但没有据此宣称所有一阶形变均可在任意全局问题中延拓。
